\documentclass[11pt,a4paper]{article}

\usepackage[T1]{fontenc}
\usepackage[utf8]{inputenc}
\usepackage[spanish,es-nodecimaldot,es-noshorthands]{babel}
\usepackage{lmodern}
\usepackage{microtype}
\usepackage[margin=2.3cm]{geometry}
\usepackage{amsmath,amssymb}
\usepackage{graphicx}
\usepackage{booktabs}
\usepackage{caption}
\usepackage{subcaption}
\usepackage{fancyhdr}
\usepackage{xcolor}
\usepackage{tikz}
\usetikzlibrary{positioning,arrows.meta,fit,backgrounds,calc}
\usepackage{csquotes}
\usepackage[backend=biber,style=numeric-comp,sorting=none]{biblatex}
\addbibresource{referencias.bib}
\usepackage[hidelinks]{hyperref}
\usepackage[spanish,noabbrev]{cleveref}

\graphicspath{{figures/}}
\captionsetup{font=small,labelfont=bf}

\pagestyle{fancy}
\fancyhf{}
\fancyhead[L]{\small Segmentación Cuaterniónica QFT}
\fancyhead[R]{\small Proyecto de Investigación I --- CIMAT}
\fancyfoot[C]{\thepage}

% Comando para valores aún no confirmados
\newcommand{\pendiente}[1]{\textcolor{red}{[#1]}}

\begin{document}

\begin{center}
  {\large Centro de Investigación en Matemáticas (CIMAT)}\\[2pt]
  Maestría en Modelación y Optimización de Procesos\\[2pt]
  \small Proyecto de Investigación I \hfill Semestre 2026-II\\[14pt]
  {\LARGE\bfseries Identificación de Lesiones Dermatológicas a través de\\[3pt]
   Segmentación por Superpíxeles, Transformada\\[3pt]
   Cuaterniónica y Clasificación mediante Random Forest\par}
  \vspace{12pt}
  Maximiliano Vaca \quad y \quad Pablo Alvarado\\[2pt]
  \small\texttt{maximiliano.vaca@cimat.mx}, \texttt{pablo.alvarado@cimat.mx}\\[2pt]
  \small Metodología propuesta por: Dr.\ Alberto López Alanís\\[2pt]
  \small Código: \url{https://github.com/MaximilianoVM/qft-medical-segmentation}
\end{center}

\begin{abstract}
\noindent
Se presenta un sistema automatizado para segmentar lesiones cutáneas en imágenes
dermatoscópicas. El método divide cada imagen en $k = 85$ superpíxeles con SLIC,
representa el color RGB de cada píxel como un cuaternión puro, aplica la
Transformada Cuaterniónica de Fourier 2D (QFT) sobre cada superpíxel y resume su
espectro en cinco descriptores (magnitud media, desviación estándar, máximo,
energía y entropía). Un clasificador Random Forest decide si cada superpíxel
pertenece a la lesión, y la máscara final se reconstruye a partir de esas
decisiones. Evaluado sobre las 379 imágenes del conjunto oficial de prueba de
ISIC 2016, el sistema obtiene \RESUMENRESULTADOS. El análisis muestra que el
método funciona bien en lesiones de alto contraste y falla cuando la lesión y la
piel tienen tonos parecidos.
\end{abstract}

\noindent\textbf{Palabras clave:} segmentación de imágenes médicas, superpíxeles
SLIC, Transformada Cuaterniónica de Fourier, Random Forest, ISIC 2016.

%=====================================================================
\section{Introducción}

La detección automática de lesiones en imágenes dermatológicas debe lidiar con
la variabilidad de la pigmentación, bordes irregulares y artefactos como vello o
reflejos. Muchos métodos clásicos procesan cada canal de color por separado o
convierten la imagen a escala de grises, con lo que pierden la correlación entre
canales. Este trabajo aborda el problema con dos ideas:

\begin{enumerate}
  \item \textbf{Superpíxeles como unidad de decisión.} SLIC agrupa la imagen en
  $k$ regiones contiguas y homogéneas \cite{achanta2012slic}. Clasificar
  regiones en lugar de píxeles reduce el costo computacional y hace que cada
  decisión use información local de textura.
  \item \textbf{Color como cuaternión.} Cada píxel RGB se codifica como un
  cuaternión puro y se analiza con la QFT
  \cite{sangwine1996fourier,ell2007hypercomplex}, de modo que los tres canales se
  transforman como una sola entidad y no como tres señales independientes.
\end{enumerate}

%=====================================================================
\section{Marco teórico}

\subsection{Presegmentación en superpíxeles (SLIC)}
SLIC adapta $k$-means a un espacio de cinco dimensiones: color CIELAB $(L,a,b)$
y posición $(x,y)$. La distancia entre un píxel $i$ y un centro $j$ combina
ambas partes:
\begin{equation}
  d_{lab} = \sqrt{(L_i-L_j)^2 + (a_i-a_j)^2 + (b_i-b_j)^2}, \qquad
  d_{xy}  = \sqrt{(x_i-x_j)^2 + (y_i-y_j)^2},
\end{equation}
\begin{equation}
  D_s = \sqrt{d_{lab}^2 + \left(\frac{m}{S}\right)^2 d_{xy}^2},
\end{equation}
donde $S = \sqrt{N/k}$ es el intervalo de muestreo para una imagen de $N$ píxeles
y $m$ controla la compacidad (en este trabajo $m = 10$ y 10 iteraciones).

\subsection{Representación cuaterniónica de imágenes RGB}
Un cuaternión es $q = w + x\,\mathbf{i} + y\,\mathbf{j} + z\,\mathbf{k}$, con
$\mathbf{i}^2 = \mathbf{j}^2 = \mathbf{k}^2 = \mathbf{ijk} = -1$. Cada píxel se
codifica como un cuaternión puro (parte real nula), con los canales normalizados
a $[0,1]$:
\begin{equation}
  q(x,y) = R(x,y)\,\mathbf{i} + G(x,y)\,\mathbf{j} + B(x,y)\,\mathbf{k}.
\end{equation}

\subsection{Transformada Cuaterniónica de Fourier 2D}
Como la multiplicación de cuaterniones no es conmutativa, existen varias
definiciones de la QFT. La implementación usa la forma \emph{bilateral}
\cite{ell2007hypercomplex}, con el núcleo $\mathbf{i}$ a la izquierda sobre las
filas y el núcleo $\mathbf{j}$ a la derecha sobre las columnas. Para una región de
$M\times N$ píxeles:
\begin{equation}
  F(u,v) = \frac{1}{\sqrt{MN}} \sum_{x=0}^{M-1}\sum_{y=0}^{N-1}
  e^{-\mathbf{i}\,2\pi \frac{ux}{M}}\; q(x,y)\; e^{-\mathbf{j}\,2\pi \frac{vy}{N}}.
  \label{eq:qft}
\end{equation}
En la práctica, \cref{eq:qft} se calcula sin aritmética de cuaterniones: se
obtiene la transformada de cada canal real con FFT en cascada (primero sobre $y$
con la unidad $\mathbf{j}$, luego sobre $x$ con $\mathbf{i}$) y las cuatro
componentes resultantes de los tres canales se combinan en un único cuaternión
por canal espectral.

\subsection{Descriptores espectrales}
De cada espectro se calcula la magnitud cuaterniónica
$|F(u,v)| = \sqrt{F_w^2 + F_x^2 + F_y^2 + F_z^2}$ y, a partir de ella, un vector
de cinco características: media, desviación estándar y máximo de $|F|$, la
energía $E = \frac{1}{MN}\sum_{u,v}|F(u,v)|^2$ y la entropía de Shannon
$H = -\sum_{u,v} p(u,v)\log_2 p(u,v)$ con $p(u,v) = |F(u,v)|^2 / \sum|F|^2$.
Las características se estandarizan por \emph{z-score} entre los superpíxeles de
cada imagen.

\subsection{Etiquetado supervisado y clasificación}
Para entrenar, cada superpíxel $S_i$ recibe una etiqueta según su traslape con la
máscara de referencia (\emph{Ground Truth}, GT):
\begin{equation}
  y_i =
  \begin{cases}
    1, & \text{si } \dfrac{|S_i \cap GT|}{|S_i|} \ge \tau, \\[4pt]
    0, & \text{en otro caso,}
  \end{cases}
\end{equation}
con $\tau = 0.9$. Un Random Forest \cite{breiman2001random} aprende a asignar
estas etiquetas a partir de los cinco descriptores. En la predicción, cada píxel
toma la clase de su superpíxel, lo que produce una máscara binaria
$M_{pred}(x,y)\in\{0,1\}$ (exportada como $0$--$255$).

\subsection{Métricas de evaluación}
La máscara predicha $P$ se compara píxel a píxel con la de referencia $G$:
\begin{equation}
  \mathrm{IoU} = \frac{|P\cap G|}{|P\cup G|},\quad
  \mathrm{Dice} = \frac{2|P\cap G|}{|P|+|G|},\quad
  \mathrm{Precisión} = \frac{|P\cap G|}{|P|},\quad
  \mathrm{Recall} = \frac{|P\cap G|}{|G|}.
\end{equation}

%=====================================================================
\section{Metodología}

El proyecto sigue el esquema propuesto por el Dr.\ Alberto López Alanís
\cite{lopez2023apuntes}, organizado en tres fases (\cref{fig:pipeline}).
Los datos provienen de ISIC 2016 \cite{gutman2016skin}: 900 imágenes de
entrenamiento y 379 de prueba, cada una con su máscara de referencia.

\begin{figure}[htbp]
\centering
\begin{tikzpicture}[
  font=\footnotesize,
  caja/.style={draw, rounded corners=2pt, fill=white, align=center,
               text width=3.4cm, minimum height=0.95cm},
  fase/.style={font=\footnotesize\bfseries, align=center},
  flecha/.style={-{Stealth[length=2mm]}, thick},
  node distance=0.42cm
]
  % Fase 1
  \node[fase] (f1) {01 Extracción de\\características};
  \node[caja, below=0.3cm of f1] (a1) {\textbf{Carga de imágenes}\\RGB + máscara GT};
  \node[caja, below=of a1] (a2) {\textbf{Superpíxeles SLIC}\\$k = 85$ regiones};
  \node[caja, below=of a2] (a3) {\textbf{QFT 2D}\\por superpíxel};
  \node[caja, below=of a3] (a4) {\textbf{Rasgos espectrales}\\magnitud, energía,\\entropía};
  % Fase 2
  \node[fase] (f2) at ($(f1)+(4.9,0)$) {02 Clasificación\\supervisada};
  \node[caja, below=0.3cm of f2] (b1) {\textbf{Normalización}\\z-score por imagen};
  \node[caja, below=of b1] (b2) {\textbf{Etiquetado GT}\\traslape $\ge \tau = 0.9$};
  \node[caja, below=of b2] (b3) {\textbf{Random Forest}\\entrenamiento};
  \node[caja, below=of b3] (b4) {\textbf{Predicción regional}\\lesión / piel sana};
  % Fase 3
  \node[fase] (f3) at ($(f2)+(4.9,0)$) {03 Reconstrucción\\y evaluación};
  \node[caja, below=0.3cm of f3] (c1) {\textbf{Reconstrucción}\\etiqueta por píxel};
  \node[caja, below=of c1] (c2) {\textbf{Máscara binaria}\\$\{0,1\}$ $\to$ $0$--$255$};
  \node[caja, below=of c2] (c3) {\textbf{Evaluación}\\IoU, Dice,\\precisión, recall};

  \begin{scope}[on background layer]
    \node[fill=blue!6, rounded corners, fit=(f1)(a4), inner sep=5pt] {};
    \node[fill=green!7, rounded corners, fit=(f2)(b4), inner sep=5pt] {};
    \node[fill=orange!8, rounded corners, fit=(f3)(c3), inner sep=5pt] {};
  \end{scope}

  \foreach \u/\v in {a1/a2, a2/a3, a3/a4, b1/b2, b2/b3, b3/b4, c1/c2, c2/c3}
    \draw[flecha] (\u) -- (\v);
  \draw[flecha] (a4.east) -- ++(0.65,0) |- (b1.west);
  \draw[flecha] (b4.east) -- ++(0.65,0) |- (c1.west);
\end{tikzpicture}
\caption{Pipeline metodológico: extracción de rasgos espectrales cuaterniónicos
por superpíxel, clasificación con Random Forest y reconstrucción y evaluación de
la máscara final.}
\label{fig:pipeline}
\end{figure}

%=====================================================================
\section{Implementación}

El código está escrito en Python y organizado en módulos dentro de \texttt{src/}:

\begin{itemize}
  \item \texttt{io\_utils.py}: carga de imágenes RGB y máscaras, y emparejamiento
        imagen--máscara.
  \item \texttt{slic\_custom.py} y \texttt{presegmentation.py}: implementación
        propia de SLIC (con alternativa \texttt{skimage.segmentation.slic}).
  \item \texttt{quaternion\_qft.py}: conversión a cuaterniones puros y QFT 2D sobre
        la caja delimitadora de cada superpíxel.
  \item \texttt{feature\_extraction.py}: descriptores espectrales y normalización.
  \item \texttt{classification.py}: etiquetado por traslape y entrenamiento de
        Random Forest, SVM y KNN.
  \item \texttt{evaluation.py}: reconstrucción de la máscara y métricas.
  \item \texttt{generate\_figures.py}: evaluación en paralelo sobre el conjunto de
        prueba y generación de todas las figuras de este reporte.
\end{itemize}

El proyecto incluye pruebas unitarias con \texttt{pytest}, un entorno
reproducible en Docker y una interfaz interactiva en Streamlit. Las figuras y
tablas de resultados se generan con \texttt{make figures}, y este documento se
compila automáticamente en GitHub Actions.

\MODELOENTRENAMIENTO

%=====================================================================
\section{Resultados y discusión}

\RESULTADOS

%=====================================================================
\section{Conclusiones y trabajo futuro}

\CONCLUSIONES

\printbibliography[title={Referencias}]

\end{document}
